\chapter{Monte Carlo and temporal-difference learning}
\label{ch:td}

Dynamic programming needs the model $p$. This chapter replaces the model by
\emph{experience}: trajectories obtained by playing. Monte Carlo methods
average complete returns; temporal-difference (TD) methods update an
estimate from the next reward and the current estimate of the next state,
which is called \emph{bootstrapping}. Their control versions, SARSA and
Q-learning, are the direct ancestors of DQN. Throughout the chapter the
values are stored in a table, one entry per state (or state-action pair);
\cref{ch:dqn} replaces the table by a neural network.

\section{Monte Carlo prediction}

Since $v_\pi(s) = \E_\pi[G_t \mid S_t = s]$, the most direct estimate of
$v_\pi(s)$ is the average of the returns observed after visits to $s$ in
episodes played with $\pi$.

\begin{algobox}{First-visit Monte Carlo prediction}
  \begin{algorithmic}[1]
    \Require policy $\pi$
    \State $V(s) \gets 0$ and $N(s) \gets 0$ for all $s$
    \For{each episode $S_0, A_0, R_1, \dots, S_T$ played with $\pi$}
      \State $G \gets 0$
      \For{$t = T-1, T-2, \dots, 0$}
        \State $G \gets R_{t+1} + \gamma\, G$
          \Comment{$G = G_t$, by \cref{prop:return-props}}
        \If{$S_t$ does not appear in $S_0, \dots, S_{t-1}$}
          \State $N(S_t) \gets N(S_t) + 1$;\quad
            $V(S_t) \gets V(S_t) + \frac{1}{N(S_t)} \left( G - V(S_t)
            \right)$
        \EndIf
      \EndFor
    \EndFor
  \end{algorithmic}
\end{algobox}

\begin{proposition}{Incremental averages}{incremental}
  \begin{enumerate}[nosep, label=(\roman*)]
    \item If $V_n = V_{n-1} + \frac{1}{n} (G_n - V_{n-1})$ with $V_0$
      arbitrary, then $V_n = \frac{1}{n} \sum_{i=1}^{n} G_i$ for $n \geq
      1$.
    \item If $V_n = V_{n-1} + \alpha (G_n - V_{n-1})$ with a constant
      $\alpha \in (0, 1]$, then
      \begin{equation}
        V_n = (1 - \alpha)^n V_0 + \sum_{i=1}^{n} \alpha (1 - \alpha)^{n-i}
        G_i,
      \end{equation}
      a weighted average in which the weight of a sample decreases
      geometrically with its age.
  \end{enumerate}
\end{proposition}

\begin{proof}
  (i) By induction: $V_1 = G_1$, and if $V_{n-1}$ is the mean of the first
  $n - 1$ samples, then $V_n = \frac{n-1}{n} V_{n-1} + \frac{1}{n} G_n =
  \frac{1}{n} \left( \sum_{i=1}^{n-1} G_i + G_n \right)$. (ii) Write
  $V_n = (1 - \alpha) V_{n-1} + \alpha G_n$ and unroll the recursion.
\end{proof}

Each first-visit return is an independent sample of $G_t$ given $S_t = s$
(episodes are independent), with mean $v_\pi(s)$ and finite variance; by
the strong law of large numbers, $V(s) \to v_\pi(s)$ almost surely for every
state visited infinitely often. Monte Carlo estimates are \emph{unbiased},
but their variance can be large: a return sums all the random rewards of
the rest of the episode. They also require complete episodes, so that
nothing is learnt before the end of a game.

\section{Temporal-difference learning}

The Bellman equation $v_\pi(s) = \E_\pi[R_{t+1} + \gamma\, v_\pi(S_{t+1})
\mid S_t = s]$ suggests replacing the return $G_t$ by the \emph{TD
target} $R_{t+1} + \gamma V(S_{t+1})$, available after a single step.

\begin{definition}{TD(0) and the TD error}{td0}
  After each transition $(S_t, A_t, R_{t+1}, S_{t+1})$ played with $\pi$,
  TD(0) updates
  \begin{equation}
    \label{eq:td0}
    V(S_t) \gets V(S_t) + \alpha\, \delta_t, \qquad
    \delta_t = R_{t+1} + \gamma\, V(S_{t+1}) - V(S_t),
  \end{equation}
  where $\alpha \in (0, 1]$ is the \emph{step size} and $\delta_t$ the
  \emph{TD error}. If $S_{t+1}$ is terminal, $V(S_{t+1}) = 0$; if the
  episode was truncated (\cref{rem:truncation}), $S_{t+1}$ is not terminal
  and its estimate is used.
\end{definition}

The TD target is \emph{biased} (it uses the current estimate $V$ instead of
$v_\pi$) but has a much smaller variance than the return: it involves one
random transition instead of a whole episode. In expectation, TD(0) moves
$V$ towards $\T^\pi V$:

\begin{proposition}{The expected TD(0) update is a contraction}{td-contraction}
  For $\alpha \in (0, 1]$, define $U_\alpha V = V + \alpha (\T^\pi V - V)$,
  the update obtained by replacing, in every state at once, the sampled
  target by its expectation. Then $U_\alpha$ is a contraction of modulus
  $1 - \alpha (1 - \gamma) < 1$ for $\norm{\cdot}_\infty$, and its unique
  fixed point is $v_\pi$.
\end{proposition}

\begin{proof}
  Given $S_t = s$, the expectation of the target is $\sum_a \pi(a \mid s)
  \sum_{s', r} p(s', r \mid s, a) [r + \gamma V(s')] = (\T^\pi V)(s)$.
  Now $U_\alpha V = (1 - \alpha) V + \alpha \T^\pi V$, so
  \begin{equation*}
    \norm{U_\alpha V - U_\alpha W}_\infty
    \leq (1 - \alpha) \norm{V - W}_\infty + \alpha \gamma \norm{V -
      W}_\infty
    = \left(1 - \alpha (1 - \gamma)\right) \norm{V - W}_\infty.
  \end{equation*}
  A fixed point satisfies $\alpha (\T^\pi V - V) = 0$, i.e.\ $V = \T^\pi
  V$, whose only solution is $v_\pi$.
\end{proof}

The real algorithm uses samples and updates one state at a time; its
convergence is a result of \emph{stochastic approximation}.

\begin{theorem}{Convergence of TD(0) (stated without proof)}{td-conv}
  Suppose every state is visited infinitely often, and the step size used
  for the $n$-th update of each state, $\alpha_n$, satisfies the
  Robbins--Monro conditions
  \begin{equation}
    \label{eq:robbins-monro}
    \sum_{n} \alpha_n = \infty, \qquad \sum_{n} \alpha_n^2 < \infty
  \end{equation}
  (for example $\alpha_n = 1/n$). Then TD(0) converges to $v_\pi$ with
  probability $1$.
\end{theorem}

The proof (Dayan, 1992; Jaakkola, Jordan and Singh, 1994; Tsitsiklis, 1994)
combines \cref{prop:td-contraction} with martingale arguments that control
the noise of the samples: the first condition of \eqref{eq:robbins-monro}
lets the estimates travel arbitrarily far, the second makes the noise
average out. In deep RL the step size is usually kept constant, and the
estimates then fluctuate around the solution instead of converging.

\section{$n$-step returns and $\lambda$-returns}

Monte Carlo and TD(0) are two extremes: the target uses all the rewards of
the episode, or only the first one followed by an estimate. Between them,
\begin{definition}{$n$-step return}{nstep}
  The \emph{$n$-step return} from time $t$ is
  \begin{equation}
    G_{t:t+n} = R_{t+1} + \gamma R_{t+2} + \dots + \gamma^{n-1} R_{t+n} +
    \gamma^n V(S_{t+n}),
  \end{equation}
  with $G_{t:t+n} = G_t$ if the episode ends before $t + n$.
\end{definition}

\begin{proposition}{Error reduction of $n$-step returns}{nstep-error}
  For a fixed estimate $V$,
  \begin{equation}
    \max_s \abs{\E_\pi\left[G_{t:t+n} \given S_t = s\right] - v_\pi(s)}
    \leq \gamma^n \max_s \abs{V(s) - v_\pi(s)}.
  \end{equation}
\end{proposition}

\begin{proof}
  $\E_\pi[G_{t:t+n} \mid S_t = s] = ((\T^\pi)^n V)(s)$: this is the
  Bellman expectation operator applied $n$ times (by induction on $n$, as in
  the proof of \cref{prop:v-q}; a terminal state contributes $0$ in both
  expressions). Since $v_\pi = (\T^\pi)^n v_\pi$ and $\T^\pi$ is a
  $\gamma$-contraction, the difference is at most $\gamma^n \norm{V -
  v_\pi}_\infty$.
\end{proof}

The larger $n$, the less the target depends on the (possibly wrong)
estimate $V$, but the more random rewards it contains: $n$ trades bias for
variance. The \emph{$\lambda$-return} averages all the $n$-step returns
with geometric weights,
\begin{equation}
  \label{eq:lambda-return}
  G_t^\lambda = (1 - \lambda) \sum_{n=1}^{\infty} \lambda^{n-1}
  G_{t:t+n}, \qquad \lambda \in [0, 1],
\end{equation}
with $G_t^0 = G_{t:t+1}$ (TD(0)) and $G_t^1 = G_t$ (Monte Carlo). We will
meet it again as the generalised advantage estimator of \cref{ch:pg}.

\begin{lemma}{The $\lambda$-return error is a sum of TD errors}{lambda-td}
  For an episode ending at time $T$, a fixed estimate $V$ (with $V = 0$ at
  the terminal state) and $\delta_t = R_{t+1} + \gamma V(S_{t+1}) -
  V(S_t)$,
  \begin{equation}
    G_t^\lambda - V(S_t) = \sum_{k=t}^{T-1} (\gamma \lambda)^{k-t}
    \delta_k.
  \end{equation}
\end{lemma}

\begin{proof}
  The $n$-step return telescopes: $G_{t:t+n} - V(S_t) = \sum_{l=0}^{n-1}
  \gamma^l \delta_{t+l}$ (the value terms cancel in pairs; beyond $T$ the
  TD errors are $0$). Substituting into \eqref{eq:lambda-return},
  \begin{equation*}
    G_t^\lambda - V(S_t) = (1 - \lambda) \sum_{n \geq 1} \lambda^{n-1}
    \sum_{l=0}^{n-1} \gamma^l \delta_{t+l}
    = \sum_{l \geq 0} \gamma^l \delta_{t+l}\, (1 - \lambda) \sum_{n > l}
    \lambda^{n-1} = \sum_{l \geq 0} (\gamma \lambda)^l \delta_{t+l},
  \end{equation*}
  for $\lambda < 1$ (the sums are finite); for $\lambda = 1$ the identity is
  the telescoping of the Monte Carlo return.
\end{proof}

\subsection{Eligibility traces}

The $\lambda$-return needs the future of the episode (the \emph{forward
view}). The lemma suggests an equivalent \emph{backward view}: each TD
error $\delta_t$ is credited to the past states, with a weight decaying by
$\gamma \lambda$ per step. The \emph{eligibility trace} $e_t(s)$ keeps
track of these weights.

\begin{algobox}{Tabular TD($\lambda$) with accumulating traces}
  \begin{algorithmic}[1]
    \For{each episode}
      \State $e(s) \gets 0$ for all $s$
      \For{each step $S_t \to S_{t+1}$ with reward $R_{t+1}$}
        \State $\delta \gets R_{t+1} + \gamma V(S_{t+1}) - V(S_t)$
        \State $e(S_t) \gets e(S_t) + 1$
        \For{all $s$}: $V(s) \gets V(s) + \alpha\, \delta\, e(s)$;\quad
          $e(s) \gets \gamma \lambda\, e(s)$
        \EndFor
      \EndFor
    \EndFor
  \end{algorithmic}
\end{algobox}

\begin{proposition}{Forward and backward views}{traces}
  If the updates of TD($\lambda$) are accumulated during an episode and
  applied at its end (offline updating), the total update of each state is
  that of the $\lambda$-return algorithm, $\alpha \sum_t (G_t^\lambda -
  V(S_t))\, \mathbb{1}[S_t = s]$.
\end{proposition}

\begin{proof}
  The trace at time $t$ is $e_t(s) = \sum_{k \leq t} (\gamma
  \lambda)^{t-k} \mathbb{1}[S_k = s]$. Exchanging the sums,
  \begin{equation*}
    \sum_t \alpha\, \delta_t\, e_t(s) = \alpha \sum_k \mathbb{1}[S_k =
    s] \sum_{t \geq k} (\gamma \lambda)^{t-k} \delta_t = \alpha \sum_k
    \mathbb{1}[S_k = s] \left( G_k^\lambda - V(S_k) \right),
  \end{equation*}
  by \cref{lem:lambda-td}, $V$ being fixed during the episode.
\end{proof}

With online updates the equivalence is only approximate; the \emph{true
online} TD($\lambda$) of van Seijen and Sutton (2014) makes it exact. Traces
cost one memory cell per state (or per feature with function
approximation, where $\mathbf{e} \gets \gamma \lambda \mathbf{e} +
\nabla \hat v(S_t)$), and give the benefits of $\lambda$-returns without
waiting for the end of the episode.

\begin{pacman}{$n$-step returns in the repository}{nstep-repo}
  The DQN agent of the repository uses $3$-step returns, and the A2C/PPO
  agent uses the advantage version of the $\lambda$-return with
  $\lambda = 0.95$. In Pacman the important rewards are rare and large
  (eating a ghost, dying, winning): with a $1$-step target, the news of a
  death travels back one state per update; with a $3$-step target it
  reaches three moves back at once.
\end{pacman}

\section{Control: SARSA and Q-learning}

To \emph{improve} the policy without a model, we estimate action values:
by \cref{rem:opt-meaning}, acting greedily with respect to $q$ needs no
model. Two issues must be solved: exploration, and the choice of the
target.

\paragraph{Exploration.} A greedy agent may never try the action that would
reveal itself to be the best. The simplest remedy is the
\emph{$\varepsilon$-greedy} policy: with probability $1 - \varepsilon$ play
a greedy action of $Q$, otherwise an action chosen uniformly at random.

\begin{proposition}{$\varepsilon$-greedy improvement}{eps-greedy}
  Let $\pi$ be \emph{$\varepsilon$-soft} ($\pi(a \mid s) \geq \varepsilon /
  \abs{\Actions}$ for all $s, a$) and let $\pi'$ be $\varepsilon$-greedy
  with respect to $q_\pi$. Then $v_{\pi'} \geq v_\pi$.
\end{proposition}

\begin{proof}
  Let $m = \abs{\Actions}$. For every $s$,
  \begin{align*}
    \sum_a \pi'(a \mid s)\, q_\pi(s, a)
    &= \frac{\varepsilon}{m} \sum_a q_\pi(s, a)
      + (1 - \varepsilon) \max_a q_\pi(s, a) \\
    &\geq \frac{\varepsilon}{m} \sum_a q_\pi(s, a) + (1 - \varepsilon)
      \sum_a \frac{\pi(a \mid s) - \varepsilon / m}{1 - \varepsilon}\,
      q_\pi(s, a) \\
    &= \sum_a \pi(a \mid s)\, q_\pi(s, a) = v_\pi(s),
  \end{align*}
  because the weights $(\pi(a \mid s) - \varepsilon / m) / (1 -
  \varepsilon)$ are nonnegative and sum to $1$, so the weighted average is
  at most the maximum. \cref{thm:pi-improvement} concludes.
\end{proof}

\paragraph{SARSA (on-policy).} The TD update for action values, using the
action actually played next:
\begin{equation}
  Q(S_t, A_t) \gets Q(S_t, A_t) + \alpha \left[ R_{t+1} + \gamma\,
  Q(S_{t+1}, A_{t+1}) - Q(S_t, A_t) \right].
\end{equation}
Its name lists the quintuple $(S_t, A_t, R_{t+1}, S_{t+1}, A_{t+1})$ it
uses. With an $\varepsilon$-greedy policy, SARSA estimates the values of
the $\varepsilon$-greedy policy itself, exploration included.

\paragraph{Q-learning (off-policy).} Watkins' Q-learning uses the best next
action instead of the one played:
\begin{equation}
  \label{eq:q-learning}
  Q(S_t, A_t) \gets Q(S_t, A_t) + \alpha \left[ R_{t+1} + \gamma \max_{a'}
  Q(S_{t+1}, a') - Q(S_t, A_t) \right].
\end{equation}
The target is a sample of the Bellman optimality operator for action
values, $(\T^\star Q)(s, a) = \sum_{s', r} p(s', r \mid s, a) [r + \gamma
\max_{a'} Q(s', a')]$, which is a $\gamma$-contraction with fixed point
$q_\star$ (same proof as \cref{lem:contraction}, using
\eqref{eq:max-ineq}). Q-learning therefore learns $q_\star$ \emph{whatever
the policy that generates the data}, provided it explores: it is
\emph{off-policy}. This property is what allows DQN to learn from a replay
memory of old transitions (\cref{ch:dqn}).

\paragraph{Expected SARSA.} Replacing the sampled next action of SARSA by
its expectation under the policy gives
\begin{equation}
  Q(S_t, A_t) \gets Q(S_t, A_t) + \alpha \left[ R_{t+1} + \gamma \sum_{a'}
  \pi(a' \mid S_{t+1})\, Q(S_{t+1}, a') - Q(S_t, A_t) \right],
\end{equation}
which removes the variance due to the choice of $A_{t+1}$. With $\pi$
greedy with respect to $Q$, expected SARSA is Q-learning; with the
$\varepsilon$-greedy behaviour policy, it is SARSA in expectation. It can
also be used off-policy, with any target policy $\pi$.

\begin{theorem}{Convergence of Q-learning (stated without proof)}{q-conv}
  If every state-action pair is updated infinitely often and the step sizes
  of each pair satisfy \eqref{eq:robbins-monro}, tabular Q-learning
  converges to $q_\star$ with probability $1$ (Watkins and Dayan, 1992).
\end{theorem}

\begin{pacman}{Learning the reduced game from experience}{td-escape}
  We return to the reduced game of \cref{pac:dp-escape} (one ghost, reach
  the power cookie), now \emph{without} using the transition probabilities:
  the agents only sample transitions. Each episode starts from a random
  state; the agents use $\varepsilon = 0.1$ and a constant $\alpha = 0.2$
  (script \code{Course/code/td\_escape.py}). We measure the success
  probability of the \emph{greedy} policy of $Q$, computed exactly as in
  \cref{ch:dp}, averaged over all the starting states:
  \begin{center}
    \begin{tabular}{@{}lcccc@{}}
      \toprule
      Episodes & $1\,000$ & $5\,000$ & $20\,000$ & $200\,000$ \\
      \midrule
      Q-learning & $0.315$ & $0.999$ & $1.000$ & $1.000$ \\
      SARSA & $0.257$ & $0.999$ & $1.000$ & $1.000$ \\
      \bottomrule
    \end{tabular}
  \end{center}
  Both learn an optimal behaviour within a few thousand games (about
  $12$ moves each), against $0.679$ for the shortest path. Their values
  differ, as the theory predicts: after $200\,000$ episodes the mean
  distance between $\max_a Q(s, a)$ and $v_\star(s)$ is $1.2$ for
  Q-learning, which estimates $q_\star$, and $3.7$ for SARSA, which
  estimates the value of its $\varepsilon$-greedy policy: with
  probability $0.1$ that policy plays a random move, sometimes into the
  ghost, so its values are lower. Q-learning's values are not exact either:
  with a constant step size they keep fluctuating
  (\cref{thm:q-conv} requires decreasing steps).
\end{pacman}

\begin{lstlisting}[caption={The SARSA and Q-learning updates
  (\code{Course/code/td\_escape.py}, excerpt)}]
r, k2 = sample(k, a)                  # one transition of the game
if k2 < 0:                            # terminal: the target is r
    Q[k, a] += alpha * (r - Q[k, a])
    break
a2 = eps_greedy(Q, k2, eps)           # next action, actually played
if method == "q-learning":
    target = r + GAMMA * Q[k2].max()  # off-policy: best next action
else:
    target = r + GAMMA * Q[k2, a2]    # SARSA: the action played
Q[k, a] += alpha * (target - Q[k, a])
\end{lstlisting}

\section{Off-policy Monte Carlo and importance sampling}

Q-learning learns about the greedy policy while following another one
thanks to the max in its target. Monte Carlo methods have no such
shortcut: a return obtained with a \emph{behaviour} policy $b$ is a sample
of the wrong expectation. Importance sampling corrects it.

\begin{proposition}{Importance sampling of returns}{is-returns}
  Let $\pi$ (target) and $b$ (behaviour) be two policies with $b(a \mid s) >
  0$ whenever $\pi(a \mid s) > 0$, and for an episode ending at $T$ let
  \begin{equation}
    \rho_{t:T-1} = \prod_{k=t}^{T-1} \frac{\pi(A_k \mid S_k)}{b(A_k \mid
    S_k)}.
  \end{equation}
  Then $\E_b\left[\rho_{t:T-1}\, G_t \given S_t = s\right] = v_\pi(s)$.
\end{proposition}

\begin{proof}
  Given $S_t = s$, the probability of a continuation $A_t, S_{t+1}, \dots,
  S_T$ under a policy $\mu$ is $\prod_{k=t}^{T-1} \mu(A_k \mid S_k)\,
  p(S_{k+1}, R_{k+1} \mid S_k, A_k)$. The ratio of these probabilities for
  $\pi$ and $b$ is $\rho_{t:T-1}$: the transition probabilities cancel.
  Hence $\E_b[\rho_{t:T-1} G_t \mid S_t = s] = \sum_{\text{cont.}}
  P_b(\text{cont.})\, \rho_{t:T-1}\, G_t = \sum_{\text{cont.}}
  P_\pi(\text{cont.})\, G_t = v_\pi(s)$, the sum running over the
  continuations of positive probability under $b$, which include those of
  positive probability under $\pi$.
\end{proof}

Averaging $\rho_{t:T-1} G_t$ over the visits of $s$ (\emph{ordinary}
importance sampling) is unbiased, but the products of ratios can be huge,
and the variance even infinite. \emph{Weighted} importance sampling divides
by the sum of the ratios instead of the number of visits: it is biased (the
bias vanishes asymptotically) but its variance is bounded, and it is
usually preferred. The same ratios, truncated to a single step, appear in
the surrogate objective of PPO (\cref{ch:ppo}).

\section{The maximisation bias}

The target of Q-learning takes the maximum of noisy estimates, which is
biased upwards.

\begin{proposition}{Maximisation bias}{max-bias}
  Let $X_1, \dots, X_m$ be random variables with finite means. Then
  $\E\left[\max_i X_i\right] \geq \max_i \E[X_i]$, with strict inequality
  as soon as, with positive probability, the maximum is attained by
  different indices.
\end{proposition}

\begin{proof}
  For every $j$, $\max_i X_i \geq X_j$, so $\E[\max_i X_i] \geq \E[X_j]$;
  take the maximum over $j$. (The function $x \mapsto \max_i x_i$ is
  convex, and the inequality is an instance of Jensen's inequality.)
\end{proof}

If the estimates $Q(s', a')$ are unbiased but noisy, $\max_{a'} Q(s', a')$
overestimates $\max_{a'} q(s', a')$, and Q-learning propagates the
overestimation through bootstrapping. \emph{Double Q-learning} (van
Hasselt, 2010) keeps two estimates: one chooses the best action, the
other evaluates it, which removes the bias when their errors are
independent. Its deep version, double DQN, is one of the improvements of
the DQN agent of the repository (\cref{ch:dqn}).

\section{Summary}

\begin{itemize}[nosep]
  \item Monte Carlo methods average complete returns: unbiased, high
    variance, need complete episodes.
  \item TD methods bootstrap from the current estimate: biased, low
    variance, learn after every step; their expected update is a
    contraction towards the Bellman fixed point.
  \item $n$-step and $\lambda$-returns interpolate between the two.
  \item SARSA (on-policy) learns the value of the policy it follows;
    Q-learning (off-policy) learns $q_\star$ from any exploratory data,
    but overestimates through the maximisation bias.
  \item On a reduced Pacman game with $5929$ states, both learn an
    optimal behaviour in a few thousand games. The full game, with its
    $10^{29}$ states, needs generalisation: the subject of the next
    chapter.
\end{itemize}
